% ===========================================================================
\section{代入的规定}\label{sec:subst}

代入 $\varphi[t/x]$ 是用项 $t$ 替换公式 $\varphi$ 中 $x$ 的所有自由出现。
代入应当保持公式的语义：代入后的公式表达的应是“把 $x$ 所指的对象换成 $t$ 所指的对象
之后，原来的陈述”。唯一需要排除的情形是\keyword{变量捕获}：$t$ 中自由出现的变量
落入了原公式某个量词的作用域，从而被该量词约束，公式的含义随之改变。

\block{封闭项与句子}
不含自由变元的公式与项统称为\keyword{封闭的}（closed）：
\begin{itemize}[leftmargin=1.3em,itemsep=1pt]
  \item \keyword{封闭公式}（closed formula），即无自由变元的公式，称为\keyword{句子}
        （sentence）。句子的真值不依赖对变量的赋值，在一个结构中或真或假；
  \item \keyword{封闭项}（closed term），即无变元的项，在结构中解释为确定的论域元素。
\end{itemize}

\begin{definition}[可自由代入]
设 $\varphi$ 为公式，$x$ 为变元，$t$ 为项。称 $t$ 在 $\varphi$ 中\emph{可自由代入} $x$
（$t$ is free for $x$ in $\varphi$），若 $x$ 的每个被替入的自由出现位置都不处于
任何 $Qy$ 的作用域之内，其中 $y\in\operatorname{Var}(t)$ 且 $Q\in\{\forall,\exists\}$。
\end{definition}

该条件是对实际替入位置的局部条件：不要求 $t$ 的变元与 $\varphi$ 的全部约束变元
互不相同，只要求替入之后 $t$ 的变元不被捕获。若 $x$ 在 $\varphi$ 中没有自由出现，
代入不改变公式。

\begin{example}[变量捕获]\label{ex:capture}
设 $\psi=\exists x\,(x=y)$，考察 $\psi[(x+1)/y]$。
\end{example}

直接机械替换得到 $\exists x\,(x=x+1)$。$\psi$ 本来的含义是“存在与 $y$ 相等的对象”，
在论域非空的结构中恒真；替换结果却表示“存在等于自身后继的对象”——
项 $x+1$ 中自由出现的 $x$ 被 $\exists x$ 捕获。正确的做法是先对约束变元\keyword{改名}，
取公式中未出现的新变元 $u$：
\[
  \psi \;\equiv\; \exists u\,(u=y)
  \quad\Longrightarrow\quad
  \psi[(x+1)/y]=\exists u\,(u=x+1).
\]
此时 $x$ 仍然自由，语义保持为“存在与 $x+1$ 相等的对象”。

\begin{remark}
约束变元的名称不影响公式的含义：对新变元 $u$，$\exists x\,\varphi$ 与
$\exists u\,\varphi[u/x]$ 语义相同。因此先改名、后代入，可以把任意代入化为
合法代入。
\end{remark}
